\chapter[Descent of \'etale morphisms]{Descent of \'etale morphisms.\\ Application to the fundamental group}
\label{IX}
% original p. 228

\section{Recollections on \'etale morphisms}
\label{IX.1}

We shall review here the properties of \'etale morphisms
(developed in~\Ref{I}) that will be of use to us, taking this
opportunity to eliminate superfluous noetherian hypotheses from
the theory. The reader will note that even if one is interested
only in noetherian schemes, the technique of descent leads one to
introduce non-noetherian schemes
(such as~$\Spec(\widehat{A}\otimes_{A}\widehat{A})$,
where~$A$ is a noetherian local ring), and in order to be able to
apply the language of fibered categories, it is important to
define notions such as \'etale morphism etc., without introducing
any noetherian restriction into them. The reader who would be
reluctant to verify or to accept that the statements below are
true without noetherian hypotheses may content himself with
accepting them under the noetherian hypotheses of~\Ref{I},
provided he introduces these same noetherian hypotheses into the
statements of the following nos., and uses
Definition~\Ref{IX.1.1} below for the non-noetherian schemes that
arise in the arguments.

\begin{definition}
\label{IX.1.1}
Let~$f\colon X\to S$ be a morphism of preschemes, and~$x$ a
point of~$X$. We say that~$f$ is \emph{\'etale at}~$x$,
or that~$X$ is \'etale over~$S$ at~$x$, if there exist an affine
open neighborhood~$U$ of~$s=f(x)$, an affine open neighborhood~$V$
of~$x$ over~$U$, a \emph{noetherian} affine scheme~$U_{0}$,
an \emph{\'etale}~\eqref{I} and affine $U_{0}$-scheme~$V_{0}$, a
morphism~$U\to U_{0}$, and a~$U$-isomorphism
$$
V\isomto V_{0}\times_{U_{0}}U\quoi.
$$
\end{definition}

Note
% original p. 229
that when~$S$ is locally noetherian, this terminology coincides
with that of \loccit We shall likewise say that~$f$ \emph{is
\'etale}, or that~$X$ \emph{is \'etale over}~$S$, if~$f$ is
\'etale at every point~$x$ of~$X$. With these definitions, the
propositions below reduce without difficulty to the noetherian
case, where they are proved in~I~\No \Ref{I.4},
\Ref{I.5},~\Ref{I.7}. For details, the reader may
consult EGA~IV\footnote{More precisely, EGA~IV 17,~18.}.

\begin{remarks}
\label{IX.1.2}
If $f$ is \'etale at $x$, then $f$ is ``\emph{of finite
presentation at}~$x$'' (VIII~\Ref{VIII.3.5}), the local ring of~$x$
in the fiber $f^{-1}(s)$ is a \emph{finite separable extension} of
$\kres(s)$, and finally $f$ is \emph{flat at}~$x$. One can show
that the converse is true, hence that Definition~\Ref{IX.1.1} is
the same as in the case where $S$ is locally noetherian, except
that the condition ``of finite type at~$x$'' must be replaced by
``of finite presentation at~$x$''. Since this result has a
delicate proof, we did not wish to give here this definition of
the notion of \'etale morphism, which does not lend itself
directly to the proof of the properties that will follow.
\end{remarks}

Let us first note that we trivially have:
\begin{proposition}
\label{IX.1.3}
If $f\colon X\to S$ is \'etale, then every morphism $f'\colon
X'\to S'$ deduced from it by a change of base $S'\to S$ is
also \'etale.
\end{proposition}

We may therefore say that the \'etale morphisms form a
\emph{fibered subcategory} of the category of
arrows in~$\Sch$
(\cf VI~\Ref{VI.11}~a)). The object of the present
expos\'e is the study of the exactness properties of
this fibered category over~$\Sch$.
\begin{proposition}
\label{IX.1.4}
Let $f\colon X\to S$ be a morphism of preschemes. For it to be
an open immersion, it is necessary and sufficient that it be
\'etale and radicial.
\end{proposition}

\Cf I~\Ref{I.5.1}. We conclude from this that if $X$ is \'etale
over~$S$, every section of~$X$ over~$S$ is an open immersion, hence,
using~\Ref{IX.1.4} again, we find:
\begin{corollary}
\label{IX.1.5}
Let $X$ be an \'etale $S$-prescheme. Then there is a
one-to-one correspondence between the set of sections of~$X$
over~$S$, and the set of
% original p. 230
open subsets $\Gamma$ of~$X$ such that the morphism $\Gamma \to S$
induced by the structural morphism is \emph{radicial} and
\emph{surjective}.
\end{corollary}

If moreover $X$ is separated over~$S$, $\Gamma$ will be a subset
of~$X$ that is both open and closed, but this does not matter. ---
Making an obvious change of base, one can put~\Ref{IX.1.5}
in the apparently more general form:
\begin{corollary}
\label{IX.1.6}
Let $X$ and $Y$ be two $S$-preschemes, $Y$ being \'etale
over~$S$. Then the map $f \mto \Gamma_f$ which associates to
every $S$-morphism $f$ from~$X$ to $Y$ the subset of~$X\times_S Y$
underlying the graph of~$f$, is a bijection of~$\Hom_S(X,Y)$ onto
the set of open subsets of~$\Gamma$ of~$X\times_S Y$ such
that the morphism $\Gamma\to X$ induced by $\pr_1$ is
\emph{radicial} and \emph{surjective}.
\end{corollary}

\begin{proposition}
\label{IX.1.7}
Let $S_0$ be the subprescheme of~$S$ defined by a quasi-coherent
Nil-ideal, \ie such that $S_0$ has the same
underlying set as~$S$. Then the functor $X \mto X\times_S
S_0$ from the category of preschemes \'etale over~$S$
to the category of preschemes \'etale over~$S_0$ is
an equivalence of categories.
\end{proposition}

The fact that this functor is fully faithful is an
immediate consequence of~\Ref{IX.1.6}. The fact that it is
essentially surjective is contained in~I~\Ref{I.8.3}. Note
that in the preceding equivalence,
$X$ is of finite type \ie quasi-finite over~$S$, (\resp finite \ie an \'etale covering
of~$S$), if and only if $X_0$ satisfies the analogous condition
over~$S_0$; same remark for the separation condition. These facts
are immediate, and are also contained in~\Ref{IX.2.4} below.

\begin{corollary}
\label{IX.1.8}
Let $A$ be a complete noetherian local ring with residue field
$k$. Then the functor $B \mto B\otimes_A k$ is an
equivalence of the category of algebras finite and
\'etale over~$A$ with the category of algebras finite and
\'etale over~$k$, (\ie composed of a finite number of
finite separable extensions of~$k$).
\end{corollary}

\begin{proposition}
\label{IX.1.9}
For $X$ to be an \emph{\'etale covering} of~$S$ \ie finite
and \'etale over~$S$, it is necessary and sufficient that $X$ be
$S$-isomorphic to the spectrum of an Algebra $\cal{A}$ over~$S$
which is a locally free Module of finite type, and such that for
every
% original p. 231
$s \in S$, $\cal{A}_s\otimes_{\cal{O}_s}\kres(s)$ is a separable
algebra over~$\kres(s)$, hence in the present case a direct product
of finite separable extensions of~$\kres(s)$.
\end{proposition}

Finally the following result is of a less elementary nature,
being the conjunction of~I~\Ref{I.8.4} and of the
\emph{existence theorem for sheaves in algebraic geometry},
(EGA~III~5; \cf also \cite{IX.1} th\ptbl 3).

\begin{theorem}
\label{IX.1.10}
Let $S$ be the spectrum of a complete noetherian local ring, $X$ a
proper $S$-scheme, $X_0$ the fiber of~$X$ at the closed point of
$S$, (so that $X_0$ is a closed subscheme of
$X$). Then the restriction functor $X' \mto X'\times_X X_0$ is
an equivalence of the category of \'etale coverings
of~$X$ with the category of \'etale coverings of~$X_0$.
\end{theorem}

\section{Submersive and universally submersive morphisms}
\label{IX.2}

\begin{definition}
\label{IX.2.1}
A morphism $g\colon S' \to S$ of preschemes is said to be
\emph{submersive}
if it is surjective, and makes~$S$ a topological quotient space
of~$S'$ (\ie a subset $U$ of~$S$ such that $f^{-1}(U)$ is
open, is open). We say that $f$ is \emph{universally
submersive}
if for every morphism $T \to S$, the morphism $f'\colon T'=S'\times_S
T \to T$ deduced from~$f$ by change of base is submersive.
\end{definition}

It is immediate that the composite of two submersive morphisms
(\resp universally submersive morphisms) is submersive
(\resp universally submersive), and that a change of base in a
universally submersive morphism gives a universally submersive
morphism (since one has done what is needed for this). If $fg$ is
submersive (\resp universally submersive), $f$ is.
\begin{examples}
\label{IX.2.2}
a) A surjective morphism which is open, or closed, is submersive,
hence a surjective universally closed or
universally open morphism is universally submersive. For example
\emph{a surjective proper morphism is universally
submersive}. On the other hand \emph{a faithfully flat and
quasi-compact morphism is universally submersive} (VIII~\Ref{VIII.4.3}). These
will be the two most important cases for us.
\end{examples}

One
% original p. 232
can apply to a submersive or universally
submersive morphism $g\colon S' \to S$ the arguments of~VIII~\Ref{VIII.4.3};
one finds in particular:
\begin{proposition}
\label{IX.2.3}
Suppose $g\colon S' \to S$ submersive. Then the following diagram
of maps is exact:
$$
\xymatrix@C=.5cm{\Ouv (S)\ar[r] &
\Ouv (S')\ar@<2pt>[r]\ar@<-2pt>[r] & \Ouv(S''),}
$$
where $S'' = S'\times_S S'$, and where $\Ouv(X)$ denotes
the set of open subsets of the prescheme~$X$.
\end{proposition}

\begin{proposition}
\label{IX.2.4}
Let $g\colon S' \to S$ be a universally submersive morphism,
$f\colon X \to Y$ an $S$-morphism, and $f'\colon X' \to Y'$ the
$S'$-morphism deduced from it by change of base. For
$f$ to be open (\resp closed), it suffices that $f'$ be so. For
$f$ to be universally open, \resp universally
closed, \resp separated, it is necessary and sufficient that $f'$
be so. If moreover $g$ is quasi-compact, and $f$ locally of finite
type, for $f$ to be proper it is necessary and sufficient that $f'$ be so.
\end{proposition}

For this last point, one notes that if $f'$ is proper, hence
quasi-compact, then $f$ is quasi-compact (VIII~\Ref{VIII.3.3}), hence
of finite type since it is locally of finite type. On the other hand
it is separated and universally closed by what
precedes, hence it is proper.
\begin{proposition}
\label{IX.2.5}
Let $S'$ be a prescheme of finite type over the spectrum $S$ of a
complete noetherian local ring, suppose that the fiber of the closed
point $s$ of~$S$ is finite, hence the local rings in $S'$ of the
points $s'$ of this fiber are finite over~$A=\cal{O}_s$. Let~$S''$ be the
scheme sum of the spectra of the $\cal{O}_{S',s'}$ in question,
considered as a finite $S$-scheme. For $g\colon S' \to
S$ to be universally submersive, it is necessary and sufficient that the
structural morphism $S''\to S'$ be surjective.
\end{proposition}

Since there is a natural $S$-morphism $S'' \to S'$, and a
finite surjective morphism is universally submersive by~\Ref{IX.2.2},
the stated condition is sufficient. Let us therefore show that if
$S''\to S'$ is not surjective, then $g$ is not universally
submersive. Indeed, let $t$ be a point of~$S$ which is not in
% original p. 233
the image of~$S''$; there then exists an $S$-scheme $T$, spectrum of a
discrete valuation ring, whose image in $S$ is
$\{s,t\}$. Note that the image of~$S''$ in $S'$ is open, since the
morphism $S'' \to S'$ is a local isomorphism, and on the other hand
this image contains $S'_s$, and does not meet $S'_t$. It follows
that the inverse image of the latter in $T'=S'\times_S T$ is
\emph{open}, and identical to the inverse image of the closed point
of~$T$. This shows that $T'\to T$ is not submersive, hence $S'\to S$
is not universally submersive.
\begin{remark}
\label{IX.2.6}
Using the criterion IV~\Ref{IV.6.3} for a constructible subset of
a noetherian space to be open, one easily finds the following
valuative criterion for a morphism
$g\colon S'\to S$ \emph{of finite type}, with $S$ locally
noetherian, to be universally submersive: it is necessary and sufficient
that for every $S$-scheme $T$, spectrum of a discrete valuation
ring, setting $T'=S'\times_S T$, the inverse image in $T'$ of the
closed point of~$T$ be non-open.
\end{remark}

\section{Descent of morphisms of \'etale preschemes}
\label{IX.3}

\begin{proposition}
\label{IX.3.1}
Let $g\colon S'\to S$ be a \emph{surjective} morphism of
preschemes, $X$ and $Y$ two preschemes over~$S$, $X'$,
$Y'$ their inverse images over~$S'$. If $Y$ is unramified over~$S$, then the canonical map
$$
\Hom_S(X,Y) \to \Hom_{S'}(X',Y')
$$
is injective.
\end{proposition}

Indeed, by virtue of~\Ref{IX.1.6}, an $S$-morphism $f\colon X\to Y$
is known when one knows the set underlying its
graph $\Gamma$, which is a subset of~$Z=X\times_SY$. Since
$$
Z'=Z\times_SS'=X'\times_{S'}Y' \to Z
$$
is surjective, (since $S'\to S$ is), this subset $\Gamma$ is
known when one knows its inverse image in $X'\times_{S'}Y'$,
which is none other than the set underlying the graph of~$f'$. Whence
the conclusion.

A subset~$\Gamma$ of~$Z$ is the graph of an $S$-morphism $f\colon
X\to Y$ if and only if it is open, and if the morphism induced
by~$\pr_1$ from~$\Gamma$ to~$Y$ is radicial
% original p. 234
and surjective \cf \Ref{IX.1}. When the first
property is satisfied, the second one is if and
only if the inverse image~$\Gamma'$ of~$\Gamma$ in~$Z'$ satisfies
the same condition VIII~\Ref{VIII.3.1}. If one knows finally that
$Z'\to Z$ is submersive, which will be the case in particular if $S'\to
S$ is universally submersive, then~$\Gamma$ is open if and
only if~$\Gamma'$ is. Thus, the set $\Hom_S(X,Y)$ is
then in one-to-one correspondence with the set of open
subsets~$\Gamma'$ of~$Z'$ such that the projection morphism
$\pr_1\colon Z'\to X'$ is radicial and surjective, (\ie corresponding to an $S'$-morphism $f'\colon X'\to Y'$), and which
are saturated for the equivalence relation defined by
$Z'\to Z$, \ie whose two inverse images in $Z''=Z'\times_Z
Z'=Z\times_S S''$ (where $S''=S'\times_S S'$), by the one and the other
projection, are equal. Now these latter are the graphs
of the two $S''$-morphisms $X''\to Y''$ deduced from~$f'$ by
change of base, by the one and the other projection $S''\to S'$. We have
thus obtained:

\begin{proposition}
\label{IX.3.2}
Let $g\colon S'\to S$ be a \emph{universally submersive} morphism
of preschemes, $S''=S'\times_S S'$, $X$ and~$Y$ two
$S$-preschemes, $X'$ and~$Y'$ their inverse images over~$S'$,
and $X''$, $Y''$ their inverse images over~$S''$. If~$Y$ is \'etale
over~$S$, the following canonical diagram of maps is exact:
$$
\xymatrix@C=.5cm{ {\Hom_S(X,Y)} \ar[r] & {\Hom_{S'}(X',Y')}
\ar@<2pt>[r]\ar@<-2pt>[r] & {\Hom_{S''}(X'',Y'')}. }
$$
\end{proposition}

Taking~$X$ and~$Y$ \'etale over~$S$, one finds the following
statement, which moreover gives back~\Ref{IX.3.2} (even when
restricting to $X=S$, a case to which one can indeed always
reduce in~\Ref{IX.3.2}, by the change of base $X\to S$);

\begin{corollary}
\label{IX.3.3}
A universally submersive morphism of preschemes is a
descent morphism for the fibered category of
preschemes \'etale over others.
\end{corollary}

I do not know, moreover, whether it is necessarily an
\emph{effective} descent morphism for the fibered category in
question, even when one makes in addition the hypothesis that~$S$ is
noetherian and~$g$ of finite type, and restricts oneself to
\'etale coverings. We shall nevertheless give in the
following no. some useful criteria of effectiveness.

\begin{corollary}
\label{IX.3.4}
Let
% original p. 235
$g\colon S'\to S$ be a universally submersive morphism whose
fibers $g^{-1}(s)$ are ``geometrically connected'',
\ie for every extension $K/\kres(s)$, $g^{-1}(s)\otimes_{\kres(s)} K$ is
connected. Then~$S'$ is connected if~$S$ is. The functor from the
category of preschemes \'etale over~$S$ to the
category
of preschemes \'etale over~$S'$ defined by~$g$ is
fully faithful.
\end{corollary}

A subset of~$S'$ which is both open and closed is
saturated for the set-theoretic equivalence relation defined
by~$g$, since the fibers are connected, hence is the inverse image
of a subset of~$S$, which is necessarily open and closed
since~$g$ is submersive. If therefore~$S$ is connected, $S'$ is. This
also implies the following result: the composite~$fg$ of two
morphisms with universally connected fibers,
$f$ being universally submersive, has universally
connected fibers; if~$S_1'$ and~$S_2'$ over~$S$ have
universally connected fibers, the same holds for~$S_1'\times_S
S_2'$. In particular, under the conditions of~\Ref{IX.3.4}, $S''$ has
universally connected fibers over~$S$. Let then~$X$ and~$Y$ be
\'etale over~$S$, and let~$u'$ be an $S'$-morphism from~$X'$ to~$Y'$;
let us prove that it is compatible with the descent data (which
implies the desired conclusion thanks to~\Ref{IX.3.3}). Now
let~$u_1''$ and~$u_2''$ be the two $S''$-morphisms $X''\to Y''$
deduced from~$u'$. The subprescheme of~$S''$ of
coincidences of~$u_1''$ and~$u_2''$ is an induced open
subprescheme, closed fiber by fiber, as the inverse image of the
diagonal prescheme of~$Y''$ over~$S''$\kern1pt\footnote{note that the
fibers of~$S'$ over~$S$ are separated!}. It is therefore the inverse
image of a subset of~$S$. Since it contains the diagonal
in~$S''$, it is identical to~$S''$, whence $u_1''=u_2''$,
qed.
